\section{Main results}
\label{sec:equiv}

We organise the central results as a four-claim pipeline, each building on the previous. All results are fully mechanised in Rocq; a mapping from theorems to artefact names and file paths is provided as supplementary material.

\subsection{Claim 1: Pre-proof correspondence}

\begin{theorem}[Pre-proof]
\label{thm:pre-proof}
Successful supercompilation yields a rooted pre-proof whose edge structure corresponds to the drive, split, and fold operations.
\end{theorem}

The three step-level correspondences---driving, splitting, folding---are proved at the graph level via a bisimulation.

\subsection{Claim 2: Cyclic proof via budget trace}

\begin{theorem}[Cyclic proof]
\label{thm:cyclic-proof}
With the trace-condition check enabled, supercompilation yields a ranked cyclic proof on a budget-instrumented trace graph.
\end{theorem}

The boolean trace check guarantees cycle-progress. A budget construction lifts the base graph to pairs $(v,k)$ where progress edges consume budget. This yields a full ranking condition.

\subsection{Claim 3: CIU soundness}

\begin{theorem}[CIU soundness]
\label{thm:ciu}
The residualised program is CIU-equivalent to the source: for every closing substitution $\sigma$ and CBN value $v$, $t[\sigma]$ terminates to $v$ iff the residual does. The typed variant follows directly.
\end{theorem}

The proof uses fuel induction with the trace check preventing non-terminating self-loops; component lemmas lift CIU through all term constructors (appendix).

\subsection{Claim 4: Example validation}

\begin{proposition}[Examples]
\label{prop:examples}
Length-map fusion, append-length normalisation, append-associativity (2-pass), take-drop fusion, and vector-index normalisation are all validated by computational reflection.
\end{proposition}
